\worldchapter{Combat}{combat}
\label{ch:combat}

\begin{multicols}{2}

\section{Start of combat}

As soon as a conflict arises, time freezes and the group determines the order in which the combatants act.

\subsection{Initiative}

Each player rolls d6 according to the \textit{quickness} value. The exploding dice rule applies here as well. The dice results are added up. The competitor with the highest score starts the fight, the others follow in the order of their results.

\begin{pexample}[Example]
A rogue with quickness 4 rolls on her speed and gets 4, 5, 1 and 17. Her result is 27.
\end{pexample}

If two combatants have the same result, the \textit{Quickness} value decides first, and if this is also the same, the \textit{Deftness} value.

\subsection{Quick Reaction}

Before the combat begins, each participant performs a \textit{apprehension} check. This check symbolises the character's ability to react quickly to new threats. If successful, the competitor receives a ``Quick Reaction'', which allows them to react before their first round of combat begins. This Quick Reaction counts as a normal action (see below), but can only be used for reactions.

If this check fails, the participant does not receive their actions until the start of the first round of combat, and cannot act before then.

Once it is the player's turn, their actions are refreshed and the Quick Reaction expires.

\section{Sequence of rounds}

The combat is divided into \textit{combat rounds}. These have the following order:
\begin{itemize}
    \item Start of round
    \item The ``Player Combat Round'' is conducted for each participant in order of initiative.
        \begin{itemize}
            \item Start of the Player Combat Round
            \item The player's actions are refreshed
            \item The player performs their actions
            \item End of player combat round
        \end{itemize}
    \item End of round
\end{itemize}

The participant may perform an action for each of his available actions (see \textit{Actors and Actions}) or save the action for a reaction in his opponent's turn (see \textit{Reactions}).

Once the last participant has acted, the next \textit{combat round} begins with the first participant.

Both ``Start of Round'' and ``End of Round'' are phases in which reactions can be made. To do this, players must save actions and perceive the last \textit{actor} directly. Some effects, such as spells, can also be carried out during these phases.

\subsection{Actor and actions}

When it is a participant's turn, they are the \textit{actor}. The \textit{actor} is the participant who is actively acting and can use or withhold their actions as they wish.

At the start of the player combat round, the player's \textit{Actions} are refreshed. The number of actions a character has is determined by their character templates. The base value for every character is 2.

``Refresh'' therefore means that all actions are available again. If the participant has already used up any actions, e.g. by actions taken in the previous round, the available actions will be reset to the maximum.

Once the actions are refreshed, the participant can act in combat. To do this, he performs actions one after another, each act taking a certain number of \textit{actions}. Actions can be, among others, the following:
\begin{itemize}
    \item To \textbf{attack} with a weapon
    \item Parry \textbf{with a weapon or object}.
    \item \textbf{Reload} a weapon
    \item \textbf{Use} an object
    \item \textbf{Evade} a melee attack.
    \item \textbf{Aim} with a melee weapon or when firing a single shot.
    \item Perform any action (see below)
    \item \textbf{Hunker} or lay on the ground (The ``Hunkered'' status effect is active, see \cref{ch:wounds}).
    \item \textbf{Stand Up}
    \item \textbf{Walk} \textit{Quickness} + 1 meter (while performing another action without consuming an action, but the minimum roll is increased by 1).
    \item \textbf{Run} \textit{Quickness} + 5 meters
    \item \textbf{Rob} \textit{Quickness} / 2 + 1 meters (rounded up), the character must be \textit{Hunkered}. (see \cref{ch:wounds}: Conditions of the character)
\end{itemize}

Actions should not be performed together, but always one after the other, because of possible reactions.

\subsubsection{Arbitrary actions}

A character can also perform any action that is not on the list. In this case, the GM must decide whether the action requires one or more actions. An action that is not on the list should normally require one action. This could be anything, such as lighting a pipe, smashing two opponents' heads together, or throwing an object. The GM decides which roll is required.

\subsection{Reactions}

When an \textit{actor} acts in combat, all participants who directly perceive the \textit{actor} may react to that action.

The following conditions must be met in order to respond to an action:
\begin{itemize}
    \item The reacting participant must directly perceive the \textit{actor}, i.e. he must hear, see or otherwise take note of his action.
    \item The reacting participant still has unused \textit{actions}.
\end{itemize}

The reaction is announced and carried out immediately after the \textit{actor's} action. However, it takes place in the game time before the action. An \textit{action} can only be followed by one \textit{reaction} from a participant. Any number of players can react to the \textit{actor} if they recognise his \textit{action}. In practice, this means that the reacting player announces his reaction after the \textit{actor} has performed his action and possibly rolled the dice. This may vary from situation to situation.

If more than one player reacts to an action, the order of reaction is determined by initiative. The player with the highest initiative reacts first, followed by the other players in descending order of initiative.

Each \textit{reaction} reduces the available \textit{actions} of the reacting participant by one.

\begin{pexample}[Example]
Hagen is involved in a fight with a robber. Hagen has attacked in his combat round, but has saved an action to be able to react.
The robber's combat round begins. The robber attacks. The GM rolls four dice and scores three hits.
Hagen's player decides that Hagen should react with a \textit{Shield Parry}. He announces his reaction to the robber's attack after the GM has made the attack. He can do this because he still has one action left and is directly aware of his opponent's attack.
The reaction now takes place in the game before the robber's attack. The shield parry rule gives Hagen a cover roll of 5+ for his round shield. He rolls for each of the robber's three hits. He rolls a 5 twice, preventing two hits. The third hit hits him.
\end{pexample}

\section{Bonus dice actions}

Bonus and destiny dice can be used in combat to gain or steal actions.

To gain an additional action, a \textit{bonus die} can be subtracted. The additional action is available immediately, even for a reaction.

If a \textit{destiny die} is spent, an action can be stolen from an opponent. This is no longer available to the opponent in his current (or next, if it is not his turn) turn. The participant who spent the destiny die has the action immediately available, even as a reaction.

Spending dice for actions does not itself take an action.

\section{Sequence of an attack}

Attacks with weapons are handled exactly the same in melee and ranged combat. The only difference is that attacks with melee weapons are thrown at the \textit{Hand to Hand combat} skill, attacks with firearms are thrown at the \textit{Shooting} skill, and attacks with throwing weapons are thrown at the \textit{Throwing} skill.

An attack has the following phases:
\begin{itemize}
    \item The \textbf{Hit Roll} determines how many hits a character achieves in an attack with a weapon. Here, the dice are rolled on the respective weapon skill, and a distinction is made between \textit{critical hits} and \textit{hits}.
    \item The \textbf{Cover roll} is available to the attacked character if he has cover. Here it is possible to avert damage even before the hits hit the armor. Shields can provide cover.
    \item Converting \textit{hits} into \textit{wounds} taking into account \textit{protection}, \textit{penetration} and \textit{critical hits}.
\end{itemize}

\subsection{The hit roll}

To carry out an attack, a roll is made with the appropriate skill (shooting, melee or throwing). The number of dice is increased by the weapon's damage potential. The \textit{minimum roll} for this roll corresponds to the character's \textit{minimum roll}.

\wbif{tirakan}{}{
Here, a possible \textit{recoil penalty} must be taken into account if the character has already fired in the same combat round.
}

The number of dice is initially equal to the character's respective skill value (shooting, hand to hand combat, throwing) plus the \textit{damage potential} of the weapon.

\wbif{tirakan}{}{
The hit roll can also be modified by other circumstances. Different attack modes and firing at the wrong distance may cause the available dice to change.
}

Each success causes a \textit{hit} to the target of the attack. How the target can prevent damage is described under \textit{Wounds and Pierce} and \textit{Cover}.

\wbif{tirakan}{}{
\subsubsection{Recoil}

Automatic weapons usually cause \textit{recoil} when attacking, which makes it difficult to re-aim at a target in a directly following attack.

If an attack with a firearm is followed \textit{directly} by another attack from the same character \textit{within a combat round}, the minimum roll and the critical hit threshold are increased by 2. This malus increases for each subsequent attack in the same combat round. Thus, a third attack has a +4 malus on the minimum roll and critical hit threshold.

The recoil can be prevented if, for example, another action is inserted between two attacks in a combat round. For example, recoil does not apply with bows because a new arrow must be placed on the string between attacks.

Weapons can have a recoil compensation. This value lowers the malus per attack. Thus, the minimum roll for a subsequent attack with a weapon with recoil compensation 1 is only raised by 1. A recoil compensation of 2 ensures that recoil is no longer relevant for the weapon.

Recoil is not accounted for across combat rounds, only within a combat round.
}

\subsection{Critical hits}

\textit{Hits} caused during the hit roll become \textit{critical hits} if they reach the value 11 during the roll. This is equivalent to an \textit{exploding die} ``thrown farther'', which then shows a result of 5+ again. Changes to the character's \textit{minimum throw} are not applied here.

\wbif{tirakan}{}{
Critical hits can only be caused by melee attacks, single shot attacks, and throwing weapons, never by burst attacks.
}

If critical hits are achieved when attacking, they are announced separately from normal hits. A single shot from a bow could thus result in ``2 crits, 3 normal hits''.

Critical hits are treated as normal hits, but will always penetrate normal armour. Only armour of the type 'protection against critical hits' can protect against critical hits, all other types of armour protection cannot prevent critical hits.

If a \textit{cover} roll occurs, critical hits must be treated separately from normal hits. So the attacked person rolls twice on his cover, once for the number of critical hits, and once for the number of normal hits.

\begin{pexample}[Example]
The mercenary Maragas rolls 4, 5, 5 and 14, giving him 2 normal hits and a critical hit. The critical hit penetrates the armour, the normal hits can be reduced by the protection of the person attacked.
\end{pexample}

\subsubsection{Megacritical hits}

If \textit{critical hits} occur, the \textit{exploding dice} can be rolled further than 11. The roll continues until no 6 is reached on the respective die.

If a die reaches a 5 again after the second roll, it is a \textit{megacritical hit}. These hits are treated as critical hits, but cause an additional wound if not prevented.

For each roll of a 5+, the number of wounds is increased. So one megacritical hit can cause a lot of wounds. The rule of 5+ results in the following limits for wounds:
\begin{itemize}
    \item \textbf{Roll 5+}: normal hit.
    \item \textbf{Roll 11+}: critical hit - ignores armor
    \item \textbf{Roll 17+}: megacritical hit - ignores armor, +1 wound
    \item \textbf{Roll 23+}: megacritical hit - ignores armor, +2 wounds
    \item \textbf{Roll 29+}: megacritic hit - ignores armor, +3 wounds
\end{itemize}
And so on.

\subsubsection{Aiming}

\wbif{tirakan}{
It is possible to aim with ranged weapons as well as with melee weapons.
}{
With melee weapons and single shot, it is possible to aim the weapon. This is not possible in burst mode.
}

The character can invest actions to aim at his target more precisely. For every 1 action, the critical hit limit is reduced by 2 for the next attack. This bonus to critical hits may not exceed the character's perception value.

If the aiming character is hit while aiming, the accumulated aiming bonus is removed.

\wbif{tirakan}{}{
\subsection{Attack modes}

The \textit{attack modes} with which the bearer of the weapon can use it are indicated with each weapon. The player chooses arbitrarily from the available modes for each attack. Switching the fire mode on modern weapons requires no action.

\subsubsection{Hand-to-hand combat}

All melee weapons have this attack mode exclusively. The character strikes with the weapon in hand-to-hand combat.
\begin{itemize}
    \item The attack can be \textit{parried}.
    \item The attack can be \textit{dodged}.
    \item The attack can cause \textit{critical hits}.
    \item For the attack, the character can \textit{aim} beforehand.
\end{itemize}

\subsubsection{Single shot}

One shot is fired per action. This applies to many modern weapons, but also to bows, slingshots and crossbows.
\begin{itemize}
    \item The attack consumes 1 ammunition.
    \item The attack \textbf{cannot} be \textit{parried}.
    \item The attack \textbf{cannot} be \textit{dodged}.
    \item The attack can cause \textit{critical hits}.
    \item The character can \textit{aim} for the attack beforehand.
\end{itemize}

\subsubsection{Burst}

The weapon is used in burst mode, a short burst of fire is delivered, which is slightly less accurate than a single shot.
\begin{itemize}
    \item 2 dice are added to the attack roll.
    \item The attack consumes 3 ammunition.
    \item The attack \textbf{cannot} be \textit{parried}.
    \item The attack \textbf{cannot} be \textit{dodged}.
    \item The attack cannot \textbf{cause} \textit{critical hits}.
    \item For the attack, the character \textbf{cannot} \textit{aim}.
\end{itemize}
}

\subsection{Incorrect distance}

Each weapon has a specified distance at which it is effective. If the target's distance differs from that specified with the weapon, there is a penalty to the hit rolls.

If the real shooting distance is less than the specified distance of the weapon, the attack is performed normally. If the distance is increased up to the double of the weapon, the minimum roll of the hit roll is increased by 2.

If the target's distance is more than twice the weapon's range away, it is not possible to shoot or attack at the target.

\subsection{Cover}

If parts of the person being attacked are hidden from the attacker's view, the rule of cover applies. It depends on how much the attacked is hidden. The cover is classified into 3 levels:
\begin{itemize}
    \item 4+ Cover: Most of the person being attacked is hidden.
    \item 5+ cover: The target is half hidden
    \item 6+ Cover: It is a bit harder to hit the target behind light cover. This effect is achieved among other things by the ``hunkered'' condition.
\end{itemize}

If the attacked has at least 6+ cover, he is allowed a cover roll after the \textit{hit roll}. For this, he rolls as many dice as the attacker had \textit{hits}. For each success (on the minimum roll according to the cover), one hit is removed.

If the attacker has scored \textit{critical hits}, the Cover Roll must be made separately for critical and normal hits to determine which hits were prevented.

\subsubsection{Shields}

Shields can be used when the character wields a one-handed weapon.

Shields can be used in two different ways.
\begin{itemize}
    \item For \textbf{Shield Block}, the shield is readied in its own turn with two actions. In subsequent combat rounds, the shield provides the cover listed below for all attacks against the character. While the shield block is active, the character's movement range is halved. The \textbf{Shield Block} is active until the character cancels it, that is, lowers the shield.
    \item The \textbf{Shield Parry} can be used spontaneously as a \textit{reaction}. It provides the below cover roll for a single attack and costs one action.
\end{itemize}

Unlike other armour, shields have a special value, the cover value. This is expressed in the form X+, meaning that shields provide this amount of cover. A round shield provides 5+ cover, so after an attack, the attacked can roll 5+ for each hit to avoid it \textit{before} the application of \textit{Protection} and \textit{Wounds}. This is possible with both \textit{Shield Parry} and \textit{Shield Block}.

\subsection{Protection and Piercing}

Any success of the \textit{hit roll} which was not prevented by \textit{cover} is a \textit{hit} on the target of the attack. Other circumstances can also cause \textit{hits}, for example an explosion can cause ``3 hits with 2 wounds each''. Here, hits can be prevented by cover.

When a character takes a \textit{hit}, they can use \textit{protection} to avoid that hit. The character has a \textit{protection pool}, which is a combination of all their armour and other effects. For each unit of protection used, one hit is prevented, possibly with additional effects (see Protection Pool).

Any hit not prevented by \textit{protection} becomes as many wounds as the weapon or effect specifies. If nothing is specified, a hit causes one wound.

\subsection{Protection Pool}

Each character has a \textit{protection pool} made up of all their armour. Each piece of armour has a certain amount of protection, which is expressed in protection units. You can find more information about armour in the \cref{ch:equipment} chapter.

When a character is attacked or otherwise hit, they can use protection from their protection pool to prevent these hits. Using protection does not cost an action, and you can use as many as you like.

The protection pool represents the armour a character wears in battle. During combat, the armour can shift, a strap can break, and a piece of armour can fall off. As a result, the pool gets smaller during the fight, which is represented by the amount of protection spent. After the battle, all the armour in the pool is restored.

The protection pool is only available during combat. When a character takes damage outside of combat, it is up to them and the GM to assess the potential damage reduction provided by armour.

\wbif{tirakan}{
\begin{pexample}[Example]
Hagen wears a linen tabard. It has a protection value of ``N N B'', meaning it provides twice the normal protection and protection against bleeding. His opponent strikes with a powerful war axe, dealing 6 normal hits and one critical hit.

Hagen's player decides to use the ``N'' protection to remove two of the normal hits. He removes the ``N'' protection from his protection pool for this fight and takes four more hits. He cannot prevent the critical hits, as neither ``N'' (normal protection) nor ``B'' (protection against bleeding) can prevent them.
\end{pexample}
}{
\begin{pexample}[Example]
John has just escaped from hospital. He is wearing a surgical gown with ``B'' protection, which protects against bleeding. He is attacked with a knife in the street. The attacker rolls ``1 critical hit, 1 wound''. John cannot prevent the critical hit, but he uses the ``B'' protection to avoid the normal hit. He removes the ``B'' protection from his protection pool for this fight and takes one wound from the critical hit.
\end{pexample}
}

\subsection{Wounds}

A \textit{wound} is added directly to the wounds taken by the character. It can only be prevented if a template, equipment or other explicitly contains a rule that modifies wounds.

\subsection{Weaponless melee}

If the character attacks without a weapon, the player rolls hit dice equal to his \textit{hand-to-hand combat} value. The minimum roll is equal to the character's minimum roll, which is usually 5+.

If the character's \textit{Strength} value is higher than 2, the \textit{Melee} melee attack has \textit{Piercing} 1.

If the character's \textit{Quickness} value is higher than 2, the character adds one die to the roll.

The range of an unarmed melee attack is 1 meter.

\subsection{Dodge}

The attacked character can dodge a melee attack as a reaction. This requires that the attacked character has an action available and can sense the attacker. Thus, an attack from behind cannot be dodged.

The value is equal to the dodge value of the character templates plus the average of speed and dexterity (rounded up). The load of armor and weapons reduces this value.

To dodge an attack, the character rolls a die to his value in \textit{Dodge}. The minimum roll for this is increased by the number of hits the opponent scores. If the attacked person scores at least one success, he has completely dodged the attack.

\section{Parry melee attacks}

Melee attacks can be parried if the attacked has a suitable melee weapon ready and an action left.

To do this, you make a \textit{reaction} roll as if you were attacking with a weapon. For each success on this roll, one of the attacker's hits is removed. \textit{Critical hits} can only be prevented by critical successes on the parry roll.

\section{Special Attacks}

There are a number of special attacks that a character can use to refine or change their attack.

\subsection{Accurate Attack}

In the accurate attack, the character aims longer to land a better hit. The exchange ratio here is 1 action for reducing the minimum roll by 1. The exchange can also go over turns. The minimum roll can be reduced by a maximum of the character's Perception value, but cannot go below 2. No other action can be taken during this time. After that, a normal attack is made with the changed values.

\subsection{Knockout Attack}

The Knockout Attack has only the intention of knocking an opponent out, but without inflicting any damage. The attacker must wield a blunt weapon, or at least strike with a blunt object. If the attack is successful, the opponent roll a resistance check. If he does not achieve as many successes as there are hits, he is knocked out.

The attack does not inflict any wounds. Cover and armor are taken into account as usual.

\subsection{Massive attack}

In a massive attack, the character gathers all his strength to deliver a massive blow. For each additional action the character invests in this attack, the number of dice for this attack increases by 3, up to a maximum of the character's strength value.

\subsection{Disarming attack}

With a disarming attack, the attacker tries to knock the weapon out of the opponent's hand. To do this, he must succeed in an attack on the weapon's arm, with a minimum roll raised by 2. The attacked person must roll on his strength or deftness after the attack, and achieve at least as many successes as the attacker had hits.

If the attacked fails to do so, he has been disarmed.

The disarming attack doesn't cause any wounds.

\subsection{Two-handed fighting}

If the character is particularly skilled in the use of a weapon, he can wield two weapons of the same type at the same time, i.e. ambidextrously. Two-handed fighting is only possible with one-handed weapons. Weapons that are wielded with both hands anyway (heavy axes, polearms, etc.) cannot be wielded in two-handed combat.

If a character wields two weapons of the same type at the same time, the character gets one more action per combat round. The weapon he wields with his secondary hand attacks with a minimum roll increased by 1.

\subsection{Support weapon}

If this is possible with the weapon being used (usually firearms except bows), the character can place the weapon on a suitable spot before using it. Supporting takes one action. If shooting with a supported weapon, the minimum roll is reduced by 1. It costs no action to pick up a propped weapon again.

\subsection{Coup de grâce}

A character can kill an opponent directly if the opponent is \textit{unconscious}, \textit{sleeping}, or \textit{dying}. To do this, the player rolls a normal attack roll. If this roll succeeds with at least one success, the opponent receives the status \textit{dead} with the level corresponding to the successes of the attack. If the opponent is already \textit{dying}, the level of the state is increased by the number of wounds of the attack.

If the attack fails, a sleeping victim is likely to awaken.

\section{Throwing objects}
\label{sec:throwing}

If an item, such as a throwing net, is thrown at a target, the character rolls to its \textit{throw} value. The minimum roll is equal to the character's minimum roll, usually 5+.

If the roll results in at least one success, the character has hit his target.

\subsection{Deviation}

If the roll on \textit{throwing} shows no success, then the roll has failed. In this case, a roll is made on the deviation.

First, a 2D6 is thrown to determine the direction of the deviation. The result of the throw gives the direction in the way of the ``clock'', seen by the throwing character looking at the target. A 3 thus deviates to the right of the target, as seen by the throwing character.

Then a D6 is thrown, which determines the distance of the deviation in meters.

The thrown object thus lands at the determined location.
\end{multicols}
